\documentclass[12pt,a4paper]{report}

%% ─── PACKAGES ────────────────────────────────────────────────────────────────
\usepackage[a4paper, top=2.5cm, bottom=2.5cm, left=3cm, right=2.5cm]{geometry}
\usepackage{fontenc}
\usepackage{inputenc}
\usepackage{lmodern}
\usepackage{microtype}
\usepackage{setspace}
\usepackage{titlesec}
\usepackage{fancyhdr}
\usepackage{graphicx}
\usepackage[table]{xcolor}
\usepackage{booktabs}
\usepackage{longtable}
\usepackage{array}
\usepackage{tabularx}
\usepackage{multirow}
\usepackage{amsmath}
\usepackage{amssymb}
\usepackage{hyperref}
\usepackage{url}
\usepackage{enumitem}
\usepackage{caption}
\usepackage{subcaption}
\usepackage{float}
\usepackage{tocloft}
\usepackage{tikz}
\usepackage{pgf}
\usepackage{pgfplots}
\pgfplotsset{compat=1.18}
\usetikzlibrary{shapes.geometric, arrows.meta, positioning, fit, backgrounds,
                calc, decorations.pathreplacing, shadows, matrix, chains}
\usepackage{mdframed}
\usepackage{listings}
\usepackage{textcomp}
\usepackage{parskip}
\usepackage{emptypage}
\usepackage{afterpage}

%% ─── COLOURS ─────────────────────────────────────────────────────────────────
\definecolor{vitblue}{RGB}{0,51,102}
\definecolor{vitgold}{RGB}{180,130,0}
\definecolor{accentblue}{RGB}{30,100,180}
\definecolor{lightblue}{RGB}{220,235,250}
\definecolor{darkgray}{RGB}{60,60,60}
\definecolor{midgray}{RGB}{120,120,120}
\definecolor{codebg}{RGB}{248,248,248}
\definecolor{inputcol}{RGB}{13,71,161}
\definecolor{enginecol}{RGB}{180,71,0}
\definecolor{mlcol}{RGB}{27,94,32}
\definecolor{aicol}{RGB}{49,27,146}
\definecolor{outputcol}{RGB}{183,28,28}
\definecolor{tabheader}{RGB}{0,51,102}

%% ─── HYPERREF SETUP ──────────────────────────────────────────────────────────
\hypersetup{
  colorlinks  = true,
  linkcolor   = vitblue,
  urlcolor    = accentblue,
  citecolor   = accentblue,
  pdftitle    = {Distillation Column AI},
  pdfauthor   = {VIT Bhopal University},
}

%% ─── TYPOGRAPHY ──────────────────────────────────────────────────────────────
\setstretch{1.4}
\setlength{\parskip}{6pt}
\setlength{\parindent}{0pt}

%% ─── CHAPTER / SECTION STYLES ────────────────────────────────────────────────
\titleformat{\chapter}[display]
  {\normalfont\Large\bfseries\color{vitblue}}
  {\chaptertitlename\ \thechapter}{10pt}{\Huge}
\titlespacing*{\chapter}{0pt}{20pt}{30pt}

\titleformat{\section}{\large\bfseries\color{vitblue}}{{\thesection}}{1em}{}
\titleformat{\subsection}{\normalsize\bfseries\color{accentblue}}{{\thesubsection}}{1em}{}
\titleformat{\subsubsection}{\normalsize\itshape\color{darkgray}}{{\thesubsubsection}}{1em}{}

%% ─── HEADER / FOOTER ─────────────────────────────────────────────────────────
\pagestyle{fancy}
\fancyhf{}
\fancyhead[L]{\small\textit{Distillation Column AI Report}}
\fancyhead[R]{\small\textit{B.Tech Project Exhibition 2024--2025}}
\fancyfoot[C]{\thepage}
\renewcommand{\headrulewidth}{0.4pt}
\renewcommand{\footrulewidth}{0pt}

%% ─── TABLE STYLES ────────────────────────────────────────────────────────────
\newcolumntype{L}[1]{>{\raggedright\arraybackslash}p{#1}}
\newcolumntype{C}[1]{>{\centering\arraybackslash}p{#1}}
\newcolumntype{R}[1]{>{\raggedleft\arraybackslash}p{#1}}

%% ─── TABLE HEADER COMMAND ────────────────────────────────────────────────────
\newcommand{\thead}[1]{%
  \cellcolor{tabheader}\textcolor{white}{\textbf{#1}}%
}

%% ─── STYLED BOX ──────────────────────────────────────────────────────────────
\newmdenv[
  linecolor=vitblue,
  linewidth=1.5pt,
  backgroundcolor=lightblue,
  innerleftmargin=10pt,
  innerrightmargin=10pt,
  innertopmargin=8pt,
  innerbottommargin=8pt,
  roundcorner=4pt,
]{infobox}

%% ─── CODE LISTING ────────────────────────────────────────────────────────────
\lstset{
  backgroundcolor=\color{codebg},
  basicstyle=\ttfamily\small,
  breaklines=true,
  frame=single,
  rulecolor=\color{midgray},
  numbers=left,
  numberstyle=\tiny\color{midgray},
}

%% ─── TOC FORMATTING ──────────────────────────────────────────────────────────
\renewcommand{\cftchapfont}{\bfseries}
\renewcommand{\cftsecfont}{}
\setlength{\cftbeforechapskip}{4pt}

%%═══════════════════════════════════════════════════════════════════════════════
%%  DOCUMENT
%%═══════════════════════════════════════════════════════════════════════════════
\begin{document}
\pagenumbering{roman}

%% ─── TITLE PAGE ──────────────────────────────────────────────────────────────
\begin{titlepage}
\begin{center}

\vspace*{0.5cm}
{\LARGE\bfseries Distillation Column AI}\\[6pt]
{\large Intelligent Prediction $\cdot$ Explainable Diagnostics $\cdot$ Real-Time Optimization}

\vspace{1.2cm}
{\Large\bfseries A Project Report}

\vspace{0.8cm}
{\large\itshape Submitted By:}\\[6pt]
{\large
Mausam Kar (24BAI10248)\\
Poorna Gupta (24BAI10091)\\
Krishna Narayan Singh (24BAI10230)\\
MD Tanzeel Khan (24BAI10359)\\
Raj Kumar Mahto (24BAI10167)\\
Nalawde Ashlesh Kumar (24BAI10407)\\
}

\vspace{0.8cm}
{\itshape In partial fulfillment for the award of the degree of}

\vspace{0.4cm}
{\large\bfseries BACHELOR OF TECHNOLOGY}\\[4pt]
{\normalsize in}\\[4pt]
{\large\bfseries COMPUTER SCIENCE AND ENGINEERING}\\[4pt]
{\large\bfseries (ARTIFICIAL INTELLIGENCE AND MACHINE LEARNING)}

\vspace{0.8cm}
%% ── Logo placeholder (TikZ drawn VIT logo) ───────────────────────────────────
\begin{tikzpicture}
  \draw[vitblue, line width=2pt] (0,0) circle (1.2cm);
  \node[vitblue, font=\bfseries\large] at (0,0.3) {VIT};
  \node[vitblue, font=\bfseries\normalsize] at (0,-0.3) {BHOPAL};
\end{tikzpicture}

\vspace{0.6cm}
{\large\bfseries SCHOOL OF COMPUTING SCIENCE ENGINEERING AND ARTIFICIAL\\[2pt]
INTELLIGENCE}

\vspace{0.4cm}
{\bfseries
VIT BHOPAL UNIVERSITY\\
KOTHRIKALAN, SEHORE\\
MADHYA PRADESH -- 466114
}

\vspace{0.5cm}
{\large September 2025}

\end{center}
\end{titlepage}

%% ─── CERTIFICATE ─────────────────────────────────────────────────────────────
\clearpage
\thispagestyle{empty}
\begin{center}
{\bfseries\large VIT BHOPAL UNIVERSITY,}\\
{\bfseries\large KOTHRIKALAN, SEHORE}\\
{\bfseries\large MADHYA PRADESH -- 466114}

\vspace{1cm}
{\Large\bfseries BONAFIDE CERTIFICATE}
\end{center}

\vspace{0.6cm}
Certified that this project report titled \textbf{``Distillation Column AI: Intelligent
Prediction, Explainable Diagnostics and Real-Time Optimization''} is the bonafide work
of \textbf{Mausam Kar (24BAI10248), Poorna Gupta (24BAI10091), Krishna Narayan
Singh (24BAI10230), Md Tanzeel Khan (24BAI10359), Raj Kumar Mahto (24BAI10167),
Nalawde Ashlesh Kumar (24BAI10407)} who carried out the project work under my
supervision. Certified further that to the best of my knowledge the work reported at this
time does not form part of any other project/research work based on which a degree or
award was conferred on an earlier occasion on this or any other candidate.

\vspace{1.5cm}
\begin{tabular}{p{7cm}p{7cm}}
\textbf{PROGRAM CHAIR} & \textbf{PROJECT GUIDE} \\[10pt]
Dr.\ Siddharth Singh Chouhan & Dr.\ Trapti Sharma \\
School of Computing Science & School of Computing Science \\
Engineering and Artificial & Engineering and Artificial \\
Intelligence & Intelligence \\
VIT BHOPAL UNIVERSITY & VIT BHOPAL UNIVERSITY \\
\end{tabular}

\vspace{1.5cm}
The Project Exhibition I Examination is held on: \textbf{25/09/2025}

%% ─── ACKNOWLEDGEMENT ─────────────────────────────────────────────────────────
\clearpage
\thispagestyle{plain}
\begin{center}
{\Large\bfseries ACKNOWLEDGEMENT}
\end{center}
\vspace{0.6cm}

First and foremost we would like to thank the Lord Almighty for His presence and
immense blessings throughout the project work.

We wish to express our heartfelt gratitude to \textbf{Dr.\ Siddharth Singh Chouhan},
Head of the Department, School of Computing Science Engineering and Artificial
Intelligence, for his valuable support and encouragement in carrying out this work.

We would like to thank our internal guide \textbf{Dr.\ Trapti Sharma}, for continually
guiding and actively participating in the project, giving valuable suggestions throughout
to complete the project work successfully.

We extend our thanks to all the technical and teaching staff of the School of Computing
Science Engineering and Artificial Intelligence, who extended directly or indirectly all
possible support.

Last but not least, we are deeply indebted to our parents who have been the greatest
support while we worked day and night for the project to make it a success.

%% ─── ABBREVIATIONS ───────────────────────────────────────────────────────────
\clearpage
\thispagestyle{plain}
\begin{center}
{\Large\bfseries LIST OF ABBREVIATIONS}
\end{center}
\vspace{0.4cm}

\begin{longtable}{|L{4cm}|L{10cm}|}
\hline
\thead{Abbreviation} & \thead{Description} \\
\hline
\endfirsthead
\hline
\thead{Abbreviation} & \thead{Description} \\
\hline
\endhead
\hline
\endfoot
AI       & Artificial Intelligence \\
ML       & Machine Learning \\
RF       & Random Forest \\
XGB      & XGBoost (Extreme Gradient Boosting) \\
SHAP     & SHapley Additive exPlanations \\
NLP      & Natural Language Processing \\
LLM      & Large Language Model \\
API      & Application Programming Interface \\
R\textsuperscript{2}        & Coefficient of Determination (R-squared) \\
RMSE     & Root Mean Square Error \\
MAE      & Mean Absolute Error \\
CV       & Cross-Validation \\
PCA      & Principal Component Analysis \\
DCS      & Distributed Control System \\
SCADA    & Supervisory Control and Data Acquisition \\
OPC-UA   & Open Platform Communications Unified Architecture \\
REST     & Representational State Transfer \\
GPU      & Graphics Processing Unit \\
CSV      & Comma-Separated Values \\
UI       & User Interface \\
UX       & User Experience \\
VLE      & Vapour-Liquid Equilibrium \\
\end{longtable}

%% ─── LIST OF TABLES ──────────────────────────────────────────────────────────
\clearpage
\thispagestyle{plain}
\begin{center}
{\Large\bfseries LIST OF TABLES}
\end{center}
\vspace{0.4cm}

\begin{longtable}{|C{2cm}|L{10cm}|C{2cm}|}
\hline
\thead{Table No.} & \thead{Title} & \thead{Page No.} \\
\hline
\endfirsthead
\hline
\thead{Table No.} & \thead{Title} & \thead{Page No.} \\
\hline
\endhead
\hline
\endfoot
1.1  & AI Capabilities at a Glance          & \pageref{tab:capabilities} \\
1.2  & Technology Stack Overview             & \pageref{tab:techstack} \\
2.1  & System Architecture Layers            & \pageref{tab:archlayers} \\
3.1  & Input Parameters Description          & \pageref{tab:inputparams} \\
3.2  & Engineered Feature Set (21 Features)  & \pageref{tab:features} \\
4.1  & Model Hyperparameters -- Random Forest & \pageref{tab:rf_params} \\
4.2  & Model Hyperparameters -- XGBoost      & \pageref{tab:xgb_params} \\
4.3  & Model Performance Metrics             & \pageref{tab:metrics} \\
4.4  & Training Pipeline Steps               & \pageref{tab:trainpipeline} \\
5.1  & SHAP Feature Attribution Summary      & \pageref{tab:shap} \\
5.2  & Anomaly Detection Layers              & \pageref{tab:anomaly} \\
6.1  & Optimization Recommendation Examples  & \pageref{tab:optim} \\
7.1  & Development Roadmap                   & \pageref{tab:roadmap} \\
7.2  & Contribution Areas                    & \pageref{tab:contrib} \\
\end{longtable}

%% ─── LIST OF FIGURES ─────────────────────────────────────────────────────────
\clearpage
\thispagestyle{plain}
\begin{center}
{\Large\bfseries LIST OF FIGURES}
\end{center}
\vspace{0.4cm}

\begin{longtable}{|C{2cm}|L{10cm}|C{2cm}|}
\hline
\thead{Fig No.} & \thead{Title} & \thead{Page No.} \\
\hline
\endfirsthead
\hline
\thead{Fig No.} & \thead{Title} & \thead{Page No.} \\
\hline
\endhead
\hline
\endfoot
2.1 & High-Level System Architecture Diagram     & \pageref{fig:highlevel} \\
2.2 & AI Intelligence Pipeline Flowchart          & \pageref{fig:pipeline} \\
2.3 & Backend Architecture Layers                 & \pageref{fig:backend} \\
2.4 & Data Flow Pattern                           & \pageref{fig:dataflow} \\
3.1 & Feature Engineering Process                 & \pageref{fig:featureeng} \\
4.1 & Model Training Workflow                     & \pageref{fig:training} \\
5.1 & SHAP Explainability Flow                    & \pageref{fig:shap} \\
5.2 & Anomaly Detection Architecture              & \pageref{fig:anomaly} \\
6.1 & Prediction and Optimization Workflow        & \pageref{fig:optflow} \\
6.2 & Conversational AI Chat Flow                 & \pageref{fig:chatflow} \\
\end{longtable}

%% ─── ABSTRACT ────────────────────────────────────────────────────────────────
\clearpage
\thispagestyle{plain}
\begin{center}
{\Large\bfseries ABSTRACT}
\end{center}
\vspace{0.4cm}

\textbf{[PURPOSE -- METHODOLOGY -- FINDINGS]}

Distillation Column AI is a full-stack intelligent platform for predicting, analysing,
and optimising ethanol purity in industrial distillation columns. The platform addresses
the long-standing challenges of traditional trial-and-error tuning in chemical process
engineering by integrating physics-informed machine learning with Google Gemini's
large language model capabilities into a unified, operator-facing dashboard.

The methodology follows a three-layer architecture: a \emph{Prediction Layer} employing
a physics-augmented Random Forest and XGBoost ensemble trained on more than 4{,}000
real distillation experiments; an \emph{Intelligence Layer} that invokes Google Gemini
Pro to generate plain-English diagnostics, interpret SHAP attributions, flag anomalies
in real time, and recommend optimisation actions; and an \emph{Interaction Layer}
providing a conversational AI copilot that responds to natural-language operator queries
in context-aware fashion.

Feature engineering transforms seven raw sensor inputs — pressure, temperature,
reflux flow, vapour flow, distillate flow, bottoms flow, and feed flow — into
21 physics-derived features encompassing derived ratios, interaction terms, efficiency
metrics, and thermodynamic approximations. A StandardScaler normalises these features
before inference. The best-performing model (auto-selected by test R\textsuperscript{2})
achieves R\textsuperscript{2}~$>$~0.98, RMSE~0.0155, and MAE~0.0122, corresponding
to a prediction error margin of $\pm$1.55\%.

Key findings demonstrate that the platform successfully delivers real-time purity
predictions within 50\,ms, meaningful natural-language diagnostics without manual
interpretation effort, SHAP-based per-prediction explainability, multi-layer anomaly
detection combining Isolation Forest with domain physics rules, a what-if simulator
enabling virtual parameter testing, and an inverse target-purity solver. The modular
architecture supports future extensions including time-series trend analysis,
multi-column dashboards, and OPC-UA integration with live DCS data streams.

%% ─── TABLE OF CONTENTS ───────────────────────────────────────────────────────
\clearpage
{\pagestyle{plain}
\tableofcontents
\clearpage}

%%═══════════════════════════════════════════════════════════════════════════════
%% MAIN MATTER
%%═══════════════════════════════════════════════════════════════════════════════
\pagenumbering{arabic}
\pagestyle{fancy}

%%─────────────────────────────────────────────────────────────────────────────
\chapter{Introduction}
%%─────────────────────────────────────────────────────────────────────────────

\section{Project Overview}

Distillation Column AI is a modern, full-stack intelligent platform built to predict,
diagnose, and optimise ethanol purity in industrial distillation columns. Operating far
beyond a conventional prediction tool, the platform acts as an \emph{AI copilot}
for chemical process engineers by combining physics-informed machine learning with
Google Gemini Pro's conversational reasoning capabilities — all surfaced through a
single Streamlit dashboard.

Industrial distillation is a high-stakes, energy-intensive process. Even small
deviations in operating parameters — reflux ratio, vapour flow, feed conditions —
propagate non-linearly into product purity. Traditional approaches rely on first-principles
simulation or operator intuition, both of which are slow to adapt and difficult to
explain. Distillation Column AI replaces this gap with an always-on AI layer that not
only predicts outcomes but explains \emph{why} a given prediction was made and
recommends the exact corrective action.

The platform operates on three interconnected layers, each addressing a distinct aspect
of the problem:

\begin{itemize}[leftmargin=1.5cm]
  \item \textbf{Prediction Layer:} A physics-informed Random Forest model trained on
  4{,}000+ real distillation experiments, transforming 7 raw sensor inputs into
  21 engineered features to predict ethanol mole fraction with R\textsuperscript{2}~$>$~0.98.

  \item \textbf{Intelligence Layer:} Google Gemini AI interprets every prediction in
  plain English, explains each decision using SHAP values, detects anomalies in
  real-time, and suggests concrete optimisation actions.

  \item \textbf{Interaction Layer:} A conversational AI assistant enabling operators
  to ask questions like \textit{``Why did purity drop?''} or \textit{``What should
  I adjust to reach 92\%?''} and receive instant, context-aware answers.
\end{itemize}

\section{Project Objectives}

The primary objectives of Distillation Column AI are:

\begin{enumerate}[leftmargin=1.8cm]
  \item \textbf{Accurate Prediction:} Deliver real-time ethanol purity predictions
  with sub-2\% error from live sensor readings using an ensemble of physics-informed
  ML models.

  \item \textbf{Explainability:} Provide per-prediction SHAP attribution so engineers
  understand which operating variable contributes most to the predicted outcome,
  building trust in AI recommendations.

  \item \textbf{Natural Language Diagnostics:} Generate plain-English diagnostic
  reports via Google Gemini Pro that contextualise predictions without requiring
  statistical literacy.

  \item \textbf{Anomaly Detection:} Identify sensor drift, mass-balance violations,
  and out-of-distribution inputs in real time using a multi-layer detection framework.

  \item \textbf{Optimisation Assistance:} Suggest actionable parameter adjustments —
  quantified in engineering units — to achieve target purity levels or improve
  energy efficiency.

  \item \textbf{Operator Empowerment:} Lower the barrier to AI adoption in process
  plants by surfacing all intelligence through a conversational, non-technical
  interface.
\end{enumerate}

\section{Target Users}

The platform is designed for two primary audiences. \textbf{Process engineers} use it
to monitor column performance in real time, validate operating conditions, and act on
AI-generated optimisation recommendations without deep data-science knowledge.
\textbf{Research scientists and students} use it to explore the physics of distillation
through what-if simulations, SHAP-driven feature importance analysis, and batch
comparison of historical operating snapshots.

\section{Key Features}

Table~\ref{tab:capabilities} summarises the twelve core AI capabilities delivered by
the platform.

\begin{table}[H]
\centering
\caption{AI Capabilities at a Glance}
\label{tab:capabilities}
\renewcommand{\arraystretch}{1.35}
\begin{tabularx}{\textwidth}{|C{0.5cm}|L{4.2cm}|C{2.5cm}|L{6.8cm}|}
\hline
\thead{\#} & \thead{Capability} & \thead{Powered By} & \thead{Description} \\
\hline
1  & Real-Time Purity Prediction     & RF / XGBoost          & Predicts ethanol mole fraction from 7 operating parameters in under 50\,ms \\
2  & Natural Language Diagnostics    & Google Gemini Pro     & Writes a plain-English report explaining the prediction and its causes \\
3  & SHAP Explainability             & SHAP Library          & Per-prediction waterfall chart showing each feature's contribution \\
4  & Anomaly Detection               & Isolation Forest      & Flags sensor drift, mass-balance violations, and out-of-range inputs \\
5  & Smart Alerts                    & Gemini + Statistics   & Context-aware warnings with specific remediation instructions \\
6  & What-If Simulator               & Model + Gemini        & Virtual parameter changes to see predicted impact before acting \\
7  & Optimisation Advisor            & Gemini + Sensitivity  & Actionable recommendations with quantified adjustments \\
8  & Target Purity Solver            & Inverse Model         & Specify desired purity and receive recommended operating conditions \\
9  & Energy Efficiency Advisor       & Gemini + Domain Rules & Balances purity targets against steam and energy consumption \\
10 & Conversational AI Chat          & Gemini Chat API       & Natural-language Q\&A about any aspect of column behaviour \\
11 & Trend Analysis                  & Time-Series Engine    & Identifies gradual degradation patterns across sequential predictions \\
12 & Batch Analysis                  & Parallel Inference    & Upload multiple snapshots and receive comparative AI analysis \\
\hline
\end{tabularx}
\end{table}

\section{Technology Stack}
\label{sec:techstack}

The platform is built on a modern, open-source Python ecosystem. Table~\ref{tab:techstack}
lists each technology and its role.

\begin{table}[H]
\centering
\caption{Technology Stack Overview}
\label{tab:techstack}
\renewcommand{\arraystretch}{1.3}
\begin{tabularx}{\textwidth}{|L{3.2cm}|L{3.5cm}|L{7.3cm}|}
\hline
\thead{Category} & \thead{Technology} & \thead{Purpose} \\
\hline
Language           & Python 3.8+           & Core programming language \\
Web Framework      & Streamlit             & Interactive AI dashboard UI \\
AI / LLM           & Google Gemini Pro API & Natural-language diagnostics, conversational AI, and report generation \\
Explainability     & SHAP                  & Per-prediction feature attribution — waterfall, force, and summary plots \\
ML Models          & scikit-learn, XGBoost & Random Forest \& XGBoost training, inference, and cross-validation \\
Anomaly Detection  & scikit-learn, scipy   & Isolation Forest and statistical bounds for multi-layer detection \\
Data Processing    & pandas, NumPy         & Data manipulation, feature engineering, and numerical computation \\
Visualisation      & Matplotlib, Plotly    & Feature importance charts, SHAP plots, interactive what-if graphs \\
Serialisation      & pickle                & Model, scaler, and feature-name persistence \\
Environment Mgmt   & python-dotenv         & Secure API key management \\
Training Env       & Google Colab          & Cloud-based model training with GPU acceleration \\
Deployment         & GitHub Codespaces     & Cloud development via devcontainer \\
\hline
\end{tabularx}
\end{table}

%%─────────────────────────────────────────────────────────────────────────────
\chapter{System Architecture}
%%─────────────────────────────────────────────────────────────────────────────

\section{Architecture Overview}

Distillation Column AI follows a layered, service-oriented architecture in which raw
sensor data flows upward through increasingly intelligent processing stages before
delivering actionable output to the operator. The design decouples the physics-based
feature engineering pipeline from the ML inference layer, and both from the Gemini
AI intelligence layer, enabling each component to evolve independently.

The overall system is structured around five horizontally separated zones — Input,
Physics Feature Engine, ML Prediction, AI Intelligence, and Output — with Google
Gemini's API acting as the cross-cutting intelligence fabric that links prediction
results to human-readable insight.

\begin{table}[H]
\centering
\caption{System Architecture Layers}
\label{tab:archlayers}
\renewcommand{\arraystretch}{1.35}
\begin{tabularx}{\textwidth}{|L{3.2cm}|L{4cm}|L{7cm}|}
\hline
\thead{Layer} & \thead{Components} & \thead{Responsibility} \\
\hline
Operator Input  & 7 sensor parameters                 & Accepts raw pressure, temperature, and flow readings from the column \\
Feature Engine  & 21 engineered features              & Transforms raw inputs into physics-meaningful derived ratios and interaction terms \\
ML Prediction   & Random Forest / XGBoost             & Produces purity prediction with confidence score in $<$50\,ms \\
AI Intelligence & Gemini Pro, SHAP, Isolation Forest  & Generates diagnostics, explains predictions, detects anomalies, and recommends actions \\
Output          & Dashboard, Chat, Reports            & Presents all results in operator-facing natural language and interactive charts \\
\hline
\end{tabularx}
\end{table}

\section{High-Level Architecture Diagram}
\label{fig:highlevel}

\begin{figure}[H]
\centering
\resizebox{\textwidth}{!}{
\begin{tikzpicture}[
  node distance=0.6cm and 1.2cm,
  box/.style={rectangle, rounded corners=3pt, minimum width=2.8cm,
              minimum height=0.8cm, align=center, font=\small\bfseries,
              draw, thick},
  inputbox/.style={box, fill=blue!12, draw=inputcol, text=inputcol},
  enginebox/.style={box, fill=orange!12, draw=enginecol, text=enginecol, text width=4.5cm},
  mlbox/.style={box, fill=green!12, draw=mlcol, text=mlcol},
  aibox/.style={box, fill=violet!12, draw=aicol, text=aicol},
  outbox/.style={box, fill=red!10, draw=outputcol, text=outputcol, text width=3.5cm},
  arrow/.style={-Stealth, thick},
]

%% INPUT ROW
\node[inputbox] (P)  at (0,0)    {Pressure P};
\node[inputbox] (T)  [right=0.2cm of P] {Temperature T1};
\node[inputbox] (L)  [right=0.2cm of T] {Reflux Flow L};
\node[inputbox] (V)  [right=0.2cm of L] {Vapour Flow V};
\node[inputbox] (D)  [right=0.2cm of V] {Distillate D};
\node[inputbox] (B)  [right=0.2cm of D] {Bottoms B};
\node[inputbox] (F)  [right=0.2cm of B] {Feed Flow F};

\begin{scope}[on background layer]
  \node[draw=inputcol, fill=blue!5, rounded corners=5pt,
        fit=(P)(F), inner sep=6pt, label=above:{\color{inputcol}\small\bfseries Operator Input Layer}] (inputgrp) {};
\end{scope}

%% ENGINE ROW
\node[enginebox, below=1.2cm of $(P)!0.5!(F)$] (eng1)
      {Derived Ratios (Reflux, D/F, B/F, Reboiler, Condenser, Feed Norm)};
\node[enginebox, right=0.4cm of eng1] (eng2)
      {Interaction Terms (Reflux$\times$Temp, Feed$\times$Reflux)};
\node[enginebox, right=0.4cm of eng2] (eng3)
      {Efficiency Metrics \& Thermo Approximations};

\begin{scope}[on background layer]
  \node[draw=enginecol, fill=orange!5, rounded corners=5pt,
        fit=(eng1)(eng3), inner sep=6pt,
        label=above:{\color{enginecol}\small\bfseries Physics Feature Engine --- 21 Features}] (enggrp) {};
\end{scope}

%% ML ROW
\node[mlbox, minimum width=3.5cm, below=1.2cm of $(eng1)!0.5!(eng3)$] (scaler)
      {StandardScaler (Normalise 21 Features)};
\node[mlbox, minimum width=4.5cm, right=0.6cm of scaler] (mlmodel)
      {Random Forest / XGBoost (Auto-Selected)};

\begin{scope}[on background layer]
  \node[draw=mlcol, fill=green!5, rounded corners=5pt,
        fit=(scaler)(mlmodel), inner sep=6pt,
        label=above:{\color{mlcol}\small\bfseries ML Prediction Layer}] (mlgrp) {};
\end{scope}

%% AI ROW
\node[aibox, minimum width=2.8cm, below=1.2cm of $(scaler)!0.5!(mlmodel)$]
      (gemini) {Gemini NLP\\Diagnostic Engine};
\node[aibox, minimum width=2.8cm, right=0.4cm of gemini] (shapai)
      {SHAP Engine\\Attribution};
\node[aibox, minimum width=2.8cm, right=0.4cm of shapai] (anomai)
      {Anomaly Detector};
\node[aibox, minimum width=2.8cm, right=0.4cm of anomai] (optai)
      {Optimisation\\Advisor};
\node[aibox, minimum width=2.8cm, right=0.4cm of optai] (chatai)
      {Conversational\\AI Chat};

\begin{scope}[on background layer]
  \node[draw=aicol, fill=violet!5, rounded corners=5pt,
        fit=(gemini)(chatai), inner sep=6pt,
        label=above:{\color{aicol}\small\bfseries AI Intelligence Layer (Gemini Core)}] (aigrp) {};
\end{scope}

%% OUTPUT ROW
\node[outbox, below=1.2cm of $(gemini)!0.5!(chatai)$] (out1)
      {Purity Prediction + Confidence};
\node[outbox, right=0.4cm of out1] (out2)
      {AI Diagnostic Report};
\node[outbox, right=0.4cm of out2] (out3)
      {Anomaly Alerts \& Recommendations};

\begin{scope}[on background layer]
  \node[draw=outputcol, fill=red!5, rounded corners=5pt,
        fit=(out1)(out3), inner sep=6pt,
        label=above:{\color{outputcol}\small\bfseries Intelligent Output Layer}] (outgrp) {};
\end{scope}

%% ARROWS between groups
\draw[arrow, inputcol] (inputgrp.south) -- (enggrp.north);
\draw[arrow, enginecol] (enggrp.south) -- (mlgrp.north);
\draw[arrow, mlcol] (mlgrp.south) -- (aigrp.north);
\draw[arrow, aicol] (aigrp.south) -- (outgrp.north);

\end{tikzpicture}
}
\caption{High-Level System Architecture}
\end{figure}

\section{AI Intelligence Pipeline}
\label{fig:pipeline}

The AI intelligence pipeline describes how a single prediction request propagates from
raw sensor input through feature engineering, ML inference, Gemini diagnostics, and
finally into the operator-facing output layer.

\begin{figure}[H]
\centering
\begin{tikzpicture}[
  node distance=0.5cm,
  proc/.style={rectangle, rounded corners=4pt, minimum width=3.2cm,
               minimum height=0.9cm, align=center, font=\small,
               draw=vitblue, fill=lightblue, thick},
  dec/.style={diamond, aspect=2.5, minimum width=3.2cm, minimum height=0.9cm,
              align=center, font=\footnotesize, draw=vitgold,
              fill=yellow!10, thick},
  term/.style={rectangle, rounded corners=8pt, minimum width=3.2cm,
               minimum height=0.9cm, align=center, font=\small\bfseries,
               draw=mlcol, fill=green!10, thick},
  arrow/.style={-Stealth, thick, vitblue},
]

\node[proc] (inp)  {Operator enters 7 sensor parameters};
\node[proc, below=of inp] (feat)  {Physics Feature Engine: generate 21 features};
\node[proc, below=of feat] (scale) {StandardScaler normalisation};
\node[proc, below=of scale] (infer) {ML Inference (RF / XGBoost)};
\node[dec,  below=0.6cm of infer] (anom)  {Anomaly\\detected?};
\node[proc, below=0.7cm of anom] (shap)  {SHAP attribution computation};
\node[proc, below=of shap] (gem)   {Gemini Pro: generate diagnostic report};
\node[proc, below=of gem] (optim)  {Optimisation advisor + what-if analysis};
\node[term, below=of optim] (dash)  {Render dashboard \& conversational chat};

%% side branch for anomaly
\node[proc, right=2.5cm of anom, text width=2.8cm] (alert)
     {Generate Smart Alert + remediation advice};

\draw[arrow] (inp)   -- (feat);
\draw[arrow] (feat)  -- (scale);
\draw[arrow] (scale) -- (infer);
\draw[arrow] (infer) -- (anom);
\draw[arrow] (anom)  -- node[left,font=\footnotesize]{No} (shap);
\draw[arrow] (anom)  -- node[above,font=\footnotesize]{Yes} (alert);
\draw[arrow] (alert) |- (shap);
\draw[arrow] (shap)  -- (gem);
\draw[arrow] (gem)   -- (optim);
\draw[arrow] (optim) -- (dash);

\end{tikzpicture}
\caption{AI Intelligence Pipeline Flowchart}
\end{figure}

\section{Backend Architecture}
\label{fig:backend}

\begin{figure}[H]
\centering
\begin{tikzpicture}[
  layer/.style={rectangle, draw=vitblue, fill=lightblue!60, rounded corners=3pt,
                minimum width=11cm, minimum height=0.8cm, align=center,
                font=\small\bfseries, thick},
  comp/.style={rectangle, draw=accentblue, fill=white, rounded corners=2pt,
               minimum width=2.6cm, minimum height=0.7cm, align=center,
               font=\footnotesize, drop shadow},
  arrow/.style={-Stealth, thick, vitblue},
]
\node[layer] (api)   at (0,0)    {Streamlit Application Layer (UI + API calls)};
\node[layer] (mid)   at (0,-1.4) {Model Inference Middleware (model.pkl · scaler.pkl · features\_names.pkl)};
\node[layer] (feat2) at (0,-2.8) {Physics Feature Engineering Module (21 features from 7 inputs)};
\node[layer] (gemini2) at (0,-4.2) {Google Gemini Pro API Layer (Diagnostics · Chat · Optimisation)};
\node[layer] (shap2) at (0,-5.6) {SHAP Explainability Engine (Waterfall · Force · Summary)};
\node[layer] (anom2) at (0,-7.0) {Anomaly Detection Engine (Isolation Forest · Statistical Bounds · Physics Rules)};
\node[layer] (data)  at (0,-8.4) {Data Layer (dataset\_distill.csv · pandas · NumPy)};

\draw[arrow] (api)    -- (mid);
\draw[arrow] (mid)    -- (feat2);
\draw[arrow] (feat2)  -- (gemini2);
\draw[arrow] (gemini2)-- (shap2);
\draw[arrow] (shap2)  -- (anom2);
\draw[arrow] (anom2)  -- (data);
\end{tikzpicture}
\caption{Backend Architecture Layers}
\end{figure}

\section{Data Flow Pattern}
\label{fig:dataflow}

\begin{figure}[H]
\centering
\resizebox{\textwidth}{!}{
\begin{tikzpicture}[
  node distance=0.4cm and 1.5cm,
  box/.style={rectangle, rounded corners=3pt, minimum width=2.4cm,
              minimum height=0.75cm, align=center, font=\small,
              draw=vitblue, fill=lightblue, thick},
  arrow/.style={-Stealth, thick, vitblue},
  label/.style={font=\footnotesize\itshape, midway},
]
\node[box] (op)  {Operator};
\node[box, right=2cm of op]  (ui)     {Streamlit UI};
\node[box, right=2cm of ui]  (feat3)  {Feature\\Engine};
\node[box, right=2cm of feat3] (ml)   {ML Model};
\node[box, right=2cm of ml]  (gemini3){Gemini\\API};

\node[box, below=2cm of $(ui)!0.5!(ml)$] (cache) {Result\\Cache};

\draw[arrow] (op)     -- node[label,above]{1. Input}   (ui);
\draw[arrow] (ui)     -- node[label,above]{2. Features}(feat3);
\draw[arrow] (feat3)  -- node[label,above]{3. Predict} (ml);
\draw[arrow] (ml)     -- node[label,above]{4. Result}  (gemini3);
\draw[arrow] (gemini3)-- ++(0,-1.8) -| node[label,near start,below]{5. Store} (cache);
\draw[arrow] (cache)  -| node[label,near end,below]{6. Render} (ui);
\draw[arrow] (ui)     -- node[label,right]{7. Display} (op);

\end{tikzpicture}
}
\caption{Data Flow Pattern}
\end{figure}

%%─────────────────────────────────────────────────────────────────────────────
\chapter{Input Parameters and Feature Engineering}
%%─────────────────────────────────────────────────────────────────────────────

\section{Input Parameters}

The platform accepts seven operating parameters as direct sensor readings from the
distillation column. Table~\ref{tab:inputparams} describes each parameter, its
physical interpretation, and its typical operating range.

\begin{table}[H]
\centering
\caption{Input Parameters Description}
\label{tab:inputparams}
\renewcommand{\arraystretch}{1.35}
\begin{tabularx}{\textwidth}{|C{1.5cm}|L{2.5cm}|L{4.5cm}|C{2.5cm}|L{3cm}|}
\hline
\thead{Symbol} & \thead{Parameter} & \thead{Physical Meaning} & \thead{Typical Range} & \thead{Unit} \\
\hline
P   & Pressure         & Column operating pressure; controls VLE equilibrium       & 0.5 -- 3.0 & atm \\
T1  & Temperature      & Top-stage temperature; proxy for condenser duty           & 50 -- 90   & \textdegree C \\
L   & Reflux Flow      & Internal liquid return rate; controls separation quality  & 500 -- 1200 & kmol/hr \\
V   & Vapour Flow      & Rising vapour rate; drives mass-transfer                  & 600 -- 1500 & kmol/hr \\
D   & Distillate Flow  & Distillate product withdrawal rate                        & 50 -- 300  & kmol/hr \\
B   & Bottoms Flow     & Bottoms product withdrawal rate                           & 100 -- 400 & kmol/hr \\
F   & Feed Flow        & Fresh feed entering the column                            & 200 -- 600 & kmol/hr \\
\hline
\end{tabularx}
\end{table}

\section{Feature Engineering}
\label{fig:featureeng}

Physics-informed feature engineering is the cornerstone of the platform's accuracy.
Rather than feeding raw sensor readings directly to the ML model, the feature engine
derives 21 domain-meaningful features that encode thermodynamic and mass-balance
relationships. This approach injects prior chemical-engineering knowledge into the
model, dramatically reducing the training data requirement while improving
generalisation to unseen operating conditions.

\begin{figure}[H]
\centering
\resizebox{\textwidth}{!}{
\begin{tikzpicture}[
  node distance=0.5cm,
  box/.style={rectangle, rounded corners=3pt, minimum width=4cm,
              minimum height=0.75cm, align=center, font=\small,
              draw, thick},
  raw/.style={box, fill=blue!12, draw=inputcol, text=inputcol},
  der/.style={box, fill=orange!12, draw=enginecol, text=enginecol, minimum width=5.5cm},
  arrow/.style={-Stealth, thick},
]
%% raw inputs
\node[raw] (r1) at (0,0)     {P, T1, L, V, D, B, F};
\node[font=\small\bfseries\color{inputcol}, above=0.15cm of r1] {7 Raw Sensor Inputs};

%% categories
\node[der] (d1) at (7, 1.6)  {Derived Ratios (6 features)};
\node[der] (d2) at (7, 0.5)  {Interaction Terms (3 features)};
\node[der] (d3) at (7,-0.6)  {Efficiency Metrics (3 features)};
\node[der] (d4) at (7,-1.7)  {Thermo Approximations (2 features)};
\node[der] (d5) at (7,-2.8)  {Original Inputs (7 features)};

%% arrows
\draw[arrow, inputcol] (r1.east) -- ++(1.2,0) |- (d1.west);
\draw[arrow, inputcol] (r1.east) -- ++(1.2,0) |- (d2.west);
\draw[arrow, inputcol] (r1.east) -- ++(1.2,0) |- (d3.west);
\draw[arrow, inputcol] (r1.east) -- ++(1.2,0) |- (d4.west);
\draw[arrow, inputcol] (r1.east) -- ++(1.2,0) |- (d5.west);

%% output node
\node[box, fill=green!12, draw=mlcol, text=mlcol, right=1.5cm of d3, text width=3.5cm] (out)
     {21 Features\\$\rightarrow$ StandardScaler\\$\rightarrow$ ML Model};

\draw[arrow, mlcol] (d1.east) -- ++(0.8,0) |- (out.north west);
\draw[arrow, mlcol] (d2.east) -- (out.west);
\draw[arrow, mlcol] (d3.east) -- (out.west);
\draw[arrow, mlcol] (d4.east) -- ++(0.8,0) |- (out.south west);
\draw[arrow, mlcol] (d5.east) -- ++(0.8,0) |- (out.south west);
\end{tikzpicture}
}
\caption{Feature Engineering Process from 7 Inputs to 21 Physics-Derived Features}
\end{figure}

Table~\ref{tab:features} enumerates the complete 21-feature set.

\begin{table}[H]
\centering
\caption{Engineered Feature Set (21 Features)}
\label{tab:features}
\renewcommand{\arraystretch}{1.3}
\begin{tabularx}{\textwidth}{|C{0.4cm}|L{4.2cm}|L{4.5cm}|L{5.3cm}|}
\hline
\thead{\#} & \thead{Feature Name} & \thead{Formula / Derivation} & \thead{Physical Significance} \\
\hline
1  & Reflux Ratio        & $L/V$              & Degree of internal liquid-vapour contact; primary separation driver \\
2  & Reboiler Duty Ratio & $V/F$              & Fraction of feed vaporised; heat input indicator \\
3  & Condenser Ratio     & $L/D$              & Ratio of reflux to distillate; purity vs.\ throughput trade-off \\
4  & Feed Normalised     & $F/(D+B)$          & Material-balance consistency check \\
5  & Distillate Fraction & $D/F$              & Product recovery fraction for distillate \\
6  & Bottoms Fraction    & $B/F$              & Product recovery fraction for bottoms \\
7  & Reflux--Temp Interaction  & $L \times T_1$     & Non-linear coupling of reflux and condenser temperature \\
8  & Reboiler--Temp Interaction & $V \times T_1$    & Heat input modulated by temperature \\
9  & Feed--Reflux Interaction   & $F \times L$      & Feed-induced liquid load effect \\
10 & Column Load         & $(L+V)/(2F)$       & Overall hydraulic loading of the column \\
11 & Separation Duty     & $L \times V / F$   & Combined driving force for separation \\
12 & Column Efficiency   & $D / (L+V)$        & Ratio of product output to internal flows \\
13 & T\_bottom Approx    & $T_1 + 20$         & Estimated bottom-stage temperature (first-order approximation) \\
14 & Pressure–Temp Ratio & $P / T_1$          & Thermodynamic state variable \\
15--21 & Original Inputs & P, T1, L, V, D, B, F & Raw sensor readings retained for direct influence \\
\hline
\end{tabularx}
\end{table}

%%─────────────────────────────────────────────────────────────────────────────
\chapter{Machine Learning Models and Performance}
%%─────────────────────────────────────────────────────────────────────────────

\section{Model Training Workflow}
\label{fig:training}

\begin{figure}[H]
\centering
\begin{tikzpicture}[
  node distance=0.45cm,
  step/.style={rectangle, rounded corners=3pt, minimum width=8cm, text width=7.5cm,
               minimum height=0.75cm, align=center, font=\small,
               draw=vitblue, fill=lightblue, thick},
  arrow/.style={-Stealth, thick, vitblue},
  lbl/.style={circle, draw=vitblue, fill=vitblue, text=white,
              font=\bfseries\footnotesize, minimum size=0.5cm},
]
\foreach \txt/\idx in {
  {Load dataset\_distill.csv (semicolon-delimited; 4000+ records)}/1,
  {Unit Conversion: Temperatures T1--T14 from Kelvin to Celsius}/2,
  {Data Cleaning: remove duplicates, nulls, outliers (1st--99th percentile)}/3,
  {Feature Engineering: create 21 physics-derived features}/4,
  {Dataset Balancing: undersample dominant mid-purity range (40\%)}/5,
  {Train / Validation / Test Split: 60\% -- 20\% -- 20\%}/6,
  {Scaling: StandardScaler fit on training data only}/7,
  {Model Training: XGBoost (500 trees) and Random Forest (200 trees)}/8,
  {5-Fold Cross-Validation on training set for both models}/9,
  {Auto-Select best model by highest test R\textsuperscript{2}}/10,
  {Residual Analysis: Predicted vs Actual, Q-Q Plot, Homoscedasticity}/11,
  {Export Artifacts: model.pkl · scaler.pkl · features\_names.pkl}/12
} {
  \pgfmathtruncatemacro{\y}{-(\idx-1)*0.95}
  \node[step] (s\idx) at (0,\y) {\txt};
  \node[lbl] (l\idx) at (-5.0,\y) {\idx};
  \ifnum\idx>1
    \pgfmathtruncatemacro{\prev}{\idx-1}
    \draw[arrow] (s\prev.south) -- (s\idx.north);
  \fi
}
\end{tikzpicture}
\caption{Model Training Workflow (12-Step Pipeline)}
\end{figure}

\section{Model Hyperparameters}

\subsection{Random Forest Configuration}

\begin{table}[H]
\centering
\caption{Model Hyperparameters --- Random Forest}
\label{tab:rf_params}
\renewcommand{\arraystretch}{1.3}
\begin{tabularx}{0.55\textwidth}{|L{4.5cm}|R{4cm}|}
\hline
\thead{Parameter} & \thead{Value} \\
\hline
\texttt{n\_estimators}     & 200 \\
\texttt{max\_depth}        & 20 \\
\texttt{min\_samples\_split} & 5 \\
\texttt{random\_state}     & 42 \\
\texttt{n\_jobs}           & $-1$ (all cores) \\
\hline
\end{tabularx}
\end{table}

\subsection{XGBoost Configuration}

\begin{table}[H]
\centering
\caption{Model Hyperparameters --- XGBoost}
\label{tab:xgb_params}
\renewcommand{\arraystretch}{1.3}
\begin{tabularx}{0.55\textwidth}{|L{4.5cm}|R{4cm}|}
\hline
\thead{Parameter} & \thead{Value} \\
\hline
\texttt{n\_estimators}      & 500 \\
\texttt{max\_depth}         & 6 \\
\texttt{learning\_rate}     & 0.03 \\
\texttt{subsample}          & 0.85 \\
\texttt{colsample\_bytree}  & 0.85 \\
\texttt{reg\_alpha}         & 0.1 \\
\texttt{reg\_lambda}        & 1.0 \\
\texttt{random\_state}      & 42 \\
\hline
\end{tabularx}
\end{table}

\section{Model Performance}

Both models are evaluated on a held-out 20\% test set and a 5-fold cross-validation
partition. The final deployed model is the one with the higher test R\textsuperscript{2}.
Table~\ref{tab:metrics} summarises the key performance metrics.

\begin{table}[H]
\centering
\caption{Model Performance Metrics}
\label{tab:metrics}
\renewcommand{\arraystretch}{1.35}
\begin{tabularx}{0.75\textwidth}{|L{5cm}|C{3cm}|C{3cm}|}
\hline
\thead{Metric} & \thead{Random Forest} & \thead{XGBoost} \\
\hline
R\textsuperscript{2} (Test Set)       & $> 0.98$ & $> 0.97$ \\
RMSE (Test Set)                        & 0.0155   & 0.0178   \\
MAE  (Test Set)                        & 0.0122   & 0.0141   \\
Prediction Error Margin                & $\pm$1.55\% & $\pm$1.78\% \\
5-Fold CV Mean R\textsuperscript{2}    & 0.976    & 0.971    \\
Inference Latency                      & $< 50$\,ms & $< 50$\,ms \\
\hline
\end{tabularx}
\end{table}

The R\textsuperscript{2}~$> 0.98$ result signifies that the model explains more than
98\% of the variance in ethanol purity, a highly competitive figure given the
inherently noisy nature of industrial sensor data. The error margin of $\pm$1.55\%
is well within the tolerance required for most commercial-grade process control
applications.

%%─────────────────────────────────────────────────────────────────────────────
\chapter{AI Intelligence Layer}
%%─────────────────────────────────────────────────────────────────────────────

\section{Google Gemini Integration}

Google Gemini Pro serves as the reasoning and communication backbone of the platform.
After each ML inference, the predicted purity value, the top SHAP contributors, and
relevant operating context are assembled into a structured prompt that is dispatched
to the Gemini Pro API. Gemini responds with a plain-English diagnostic report covering
four areas: what the prediction means for product quality; which parameters are driving
the result and why; whether any action is recommended; and an energy-efficiency
commentary. This report is then rendered directly on the Streamlit dashboard alongside
the numerical prediction.

The conversational chat endpoint uses Gemini's Chat API with full conversation
history, enabling follow-up questions such as \textit{``Given the last prediction,
what reflux setting would push purity above 90\%?''} The platform maintains a rolling
message buffer ensuring context continuity across multi-turn interactions.

\section{SHAP Explainability Engine}
\label{fig:shap}

SHAP (SHapley Additive exPlanations) decomposes each prediction into additive
contributions from individual features, providing a mathematically rigorous answer to
the question: \textit{``Why did the model output this value?''} The platform computes
SHAP values using a TreeExplainer — the computationally efficient variant for
ensemble tree models — for every prediction request.

\begin{figure}[H]
\centering
\resizebox{\textwidth}{!}{
\begin{tikzpicture}[
  node distance=0.5cm,
  box/.style={rectangle, rounded corners=3pt, minimum width=4.5cm, text width=4cm,
              minimum height=0.8cm, align=center, font=\small,
              draw=aicol, fill=violet!8, thick},
  arrow/.style={-Stealth, thick, aicol},
]
\node[box] (pred) {ML Prediction Result};
\node[box, below=of pred] (tree) {TreeExplainer instantiation\\(model + background dataset)};
\node[box, below=of tree] (shap3) {Compute SHAP values\\for current feature vector};
\node[box, below=of shap3] (rank) {Rank features by\\absolute SHAP magnitude};
\node[box, below left=0.8cm and 1cm of rank] (waterfall) {Waterfall Chart\\(per-prediction)};
\node[box, below right=0.8cm and 1cm of rank] (summary)  {Summary / Force Plot\\(global trends)};
\node[box, below=2cm of rank] (gemini4) {Gemini interprets top 5\\SHAP features in plain English};

\draw[arrow] (pred)     -- (tree);
\draw[arrow] (tree)     -- (shap3);
\draw[arrow] (shap3)    -- (rank);
\draw[arrow] (rank)     -| (waterfall);
\draw[arrow] (rank)     -| (summary);
\draw[arrow] (waterfall)|- (gemini4);
\draw[arrow] (summary)  |- (gemini4);
\end{tikzpicture}
}
\caption{SHAP Explainability Computation Flow}
\end{figure}

Table~\ref{tab:shap} shows a representative attribution for a typical high-purity
operating point.

\begin{table}[H]
\centering
\caption{SHAP Feature Attribution Summary (Representative Prediction)}
\label{tab:shap}
\renewcommand{\arraystretch}{1.3}
\begin{tabularx}{0.72\textwidth}{|C{0.5cm}|L{5cm}|R{3.5cm}|R{3.5cm}|}
\hline
\thead{\#} & \thead{Feature} & \thead{SHAP Value} & \thead{Direction} \\
\hline
1 & Reflux Ratio ($L/V$)           & $+0.042$ & Increases purity \\
2 & Reboiler Duty Ratio ($V/F$)    & $+0.031$ & Increases purity \\
3 & Reflux--Temp Interaction       & $+0.018$ & Increases purity \\
4 & Pressure (P)                   & $-0.011$ & Decreases purity \\
5 & Bottoms Fraction ($B/F$)       & $-0.007$ & Decreases purity \\
\hline
\end{tabularx}
\end{table}

\section{Anomaly Detection Architecture}
\label{fig:anomaly}

The anomaly detection framework combines three independent detection layers. When any
layer flags an anomaly, a smart alert is generated and Gemini is invoked to produce
a plain-English explanation and recommended corrective action.

\begin{figure}[H]
\centering
\resizebox{\textwidth}{!}{
\begin{tikzpicture}[
  node distance=0.5cm,
  layer3/.style={rectangle, rounded corners=3pt, minimum width=10cm, text width=9.5cm,
                minimum height=0.85cm, align=center, font=\small\bfseries,
                draw=outputcol, fill=red!8, thick},
  det/.style={rectangle, rounded corners=3pt, minimum width=3cm,
              minimum height=0.75cm, align=center, font=\small,
              draw=midgray, fill=white, thick},
  arrow/.style={-Stealth, thick},
]
\node[layer3] (l1) at (0,0)    {Layer 1 --- Statistical Bounds: compare each input against training percentile limits};
\node[layer3] (l2) at (0,-1.3) {Layer 2 --- Isolation Forest: unsupervised anomaly scorer (contamination=0.05)};
\node[layer3] (l3) at (0,-2.6) {Layer 3 --- Physics Rules: mass-balance check $|D + B - F| / F < 5\%$};
\node[layer3] (l4) at (0,-3.9) {Layer 4 --- Temporal Drift: slope of last-N predictions for trend detection};

\node[det] (agg)  at (0,-5.2)  {Aggregate: any layer flags $\Rightarrow$ anomaly};
\node[det] (alert2) at (-3.5,-6.5) {Smart Alert (in-app)};
\node[det] (gem5)   at (0,-6.5)    {Gemini explanation};
\node[det] (rem)    at (3.5,-6.5)  {Remediation advice};

\draw[-Stealth, thick] (l4) -- (agg);
\draw[-Stealth, thick] (agg) -| (alert2);
\draw[-Stealth, thick] (agg) -- (gem5);
\draw[-Stealth, thick] (agg) -| (rem);
\end{tikzpicture}
}
\caption{Multi-Layer Anomaly Detection Architecture}
\end{figure}

Table~\ref{tab:anomaly} summarises the four detection layers, the technique used, and
the type of anomaly each layer is designed to catch.

\begin{table}[H]
\centering
\caption{Anomaly Detection Layers}
\label{tab:anomaly}
\renewcommand{\arraystretch}{1.35}
\begin{tabularx}{\textwidth}{|C{0.5cm}|L{3.5cm}|L{3cm}|L{7cm}|}
\hline
\thead{\#} & \thead{Layer} & \thead{Technique} & \thead{Anomaly Type Detected} \\
\hline
1 & Statistical Bounds  & Percentile limits     & Input values outside the 1st--99th percentile of training data (sensor malfunction or extreme operating regime) \\
2 & Isolation Forest    & Unsupervised ML       & Multivariate outliers — combinations of parameters that are individually plausible but collectively anomalous \\
3 & Physics Rules       & Mass-balance check    & Violations of $D + B \approx F$ within 5\% — indicating flow-meter error or column flooding \\
4 & Temporal Drift      & Slope detection       & Gradual monotonic drift in sequential predictions signalling membrane fouling, packing degradation, or equipment wear \\
\hline
\end{tabularx}
\end{table}

%%─────────────────────────────────────────────────────────────────────────────
\chapter{Application Workflows}
%%─────────────────────────────────────────────────────────────────────────────

\section{Prediction and Optimisation Workflow}
\label{fig:optflow}

\begin{figure}[H]
\centering
\resizebox{\textwidth}{!}{
\begin{tikzpicture}[
  node distance=0.5cm,
  proc2/.style={rectangle, rounded corners=3pt, minimum width=5cm, text width=4.5cm,
               minimum height=0.8cm, align=center, font=\small,
               draw=mlcol, fill=green!8, thick},
  dec2/.style={diamond, aspect=2.8, minimum width=5cm,
               align=center, font=\footnotesize,
               draw=vitgold, fill=yellow!12, thick},
  arrow2/.style={-Stealth, thick, mlcol},
]
\node[proc2] (s1) {Engineer inputs 7 operating parameters};
\node[proc2, below=of s1] (s2) {Feature engine generates 21 features};
\node[proc2, below=of s2] (s3) {ML model predicts ethanol mole fraction};
\node[dec2,  below=0.6cm of s3] (q1) {Purity $\geq$ target?};

\node[proc2, right=2.5cm of q1, text width=4cm] (opt)
     {Optimisation Advisor:\\suggest parameter changes};
\node[proc2, below=0.7cm of q1] (s4) {SHAP + Gemini diagnostic report};
\node[proc2, below=of s4] (s5) {Energy efficiency analysis};
\node[proc2, below=of s5] (s6) {What-if simulator available for virtual testing};
\node[proc2, below=of s6] (s7) {Render results on dashboard};

\draw[arrow2] (s1)   -- (s2);
\draw[arrow2] (s2)   -- (s3);
\draw[arrow2] (s3)   -- (q1);
\draw[arrow2] (q1)   -- node[right,font=\footnotesize]{Yes}  (s4);
\draw[arrow2] (q1)   -- node[above,font=\footnotesize]{No}   (opt);
\draw[arrow2] (opt)  |- (s4);
\draw[arrow2] (s4)   -- (s5);
\draw[arrow2] (s5)   -- (s6);
\draw[arrow2] (s6)   -- (s7);
\end{tikzpicture}
}
\caption{Prediction and Optimisation Workflow}
\end{figure}

\section{What-If Simulator}

The what-if simulator allows engineers to test hypothetical parameter changes without
touching the actual column. The operator adjusts any of the seven input sliders in the
``Simulation'' panel, and the platform immediately reruns the inference pipeline on the
modified feature vector. The resulting purity change is compared against the baseline
prediction, and Gemini Pro generates a commentary explaining which feature changes
drove the delta. A sensitivity analysis heat-map shows the marginal impact of
incrementally varying each parameter across its operating range.

Representative output for an optimisation recommendation:
\begin{infobox}
\textbf{Optimisation Advisor:} To reach 90\% ethanol purity from the current 87.3\%,
increase reflux flow L from 780 to 842~kmol/hr ($+$8\%). This change increases the
reflux ratio from 1.30 to 1.40, placing the column in the higher-purity operating
regime identified by the SHAP analysis. Expected purity gain: $+$3.2\%.
Alternative: reduce vapour flow V by 5\% to save energy; purity drops approximately
1\% to 86.3\%. Current column efficiency score: 8.4\,/\,10.
\end{infobox}

Table~\ref{tab:optim} lists representative optimisation recommendations for three
common operating scenarios.

\begin{table}[H]
\centering
\caption{Optimisation Recommendation Examples}
\label{tab:optim}
\renewcommand{\arraystretch}{1.35}
\begin{tabularx}{\textwidth}{|L{4cm}|L{4cm}|L{6cm}|}
\hline
\thead{Scenario} & \thead{Recommended Action} & \thead{Expected Outcome} \\
\hline
Purity below 85\%    & Increase L by 10\% (780 $\to$ 858\,kmol/hr) & Purity rises to 87--89\%; reflux ratio increases from 1.30 to 1.43 \\
Energy cost too high & Reduce V by 8\% while holding L constant & Energy savings of $\approx$8\%; purity drops $\approx$1.5\%, remains above spec \\
Mass-balance anomaly & Check D and B flow meters; verify $D+B \approx F$ & Resolve sensor drift; prevents false predictions \\
\hline
\end{tabularx}
\end{table}

\section{Conversational AI Chat Workflow}
\label{fig:chatflow}

\begin{figure}[H]
\centering
\begin{tikzpicture}[
  node distance=0.5cm,
  box/.style={rectangle, rounded corners=3pt, minimum width=4.5cm,
              minimum height=0.8cm, align=center, font=\small,
              draw=aicol, fill=violet!8, thick},
  arrow/.style={-Stealth, thick, aicol},
]
\node[box] (q)   {Operator types natural-language question};
\node[box, below=of q]  (hist) {System assembles conversation history\\+ current prediction context};
\node[box, below=of hist] (prompt) {Construct structured Gemini prompt:\\context + question + SHAP summary};
\node[box, below=of prompt] (call) {Gemini Chat API call};
\node[box, below=of call] (resp) {Parse Gemini response};
\node[box, below=of resp] (disp) {Append to chat window \& store in history};

\draw[arrow] (q)    -- (hist);
\draw[arrow] (hist) -- (prompt);
\draw[arrow] (prompt)-- (call);
\draw[arrow] (call) -- (resp);
\draw[arrow] (resp) -- (disp);

%% loop-back arrow
\draw[arrow] (disp.east) -- ++(1.2,0) |- node[right,font=\footnotesize]{Follow-up\\question} (hist.east);
\end{tikzpicture}
\caption{Conversational AI Chat Flow}
\end{figure}

%%─────────────────────────────────────────────────────────────────────────────
\chapter{Project Structure and Deployment}
%%─────────────────────────────────────────────────────────────────────────────

\section{Repository Structure}

The project is organised as a single-repository Python application with clear
separation between training artifacts, the Streamlit application, and the physics
feature module:

\begin{lstlisting}
Distillation-Column-Prediction/
|-- app.py                    # Streamlit dashboard entry-point
|-- model_training_script.py  # Full training pipeline (Google Colab)
|-- feature_engineering.py    # Physics-informed feature engine
|-- anomaly_detection.py      # Multi-layer anomaly detection module
|-- gemini_client.py          # Google Gemini API wrapper
|-- model.pkl                 # Serialised ML model (auto-selected)
|-- scaler.pkl                # Fitted StandardScaler
|-- features_names.pkl        # Ordered list of 21 feature names
|-- dataset_distill.csv       # Raw training dataset (4000+ records)
|-- feature_importance.png    # Pre-computed SHAP summary plot
|-- requirements.txt          # Python dependencies
|-- .devcontainer/            # GitHub Codespaces configuration
|-- LICENSE                   # MIT licence
`-- README.md                 # Project documentation
\end{lstlisting}

\section{Quick Start}

Deploying the platform locally requires Python 3.8 or later. After cloning the
repository and installing dependencies with \texttt{pip install -r requirements.txt},
the operator places a valid Google Gemini API key in a \texttt{.env} file under the
key \texttt{GOOGLE\_API\_KEY}. The pre-trained model artifacts (\texttt{model.pkl},
\texttt{scaler.pkl}, \texttt{features\_names.pkl}) must be present in the root
directory, either downloaded from the repository or generated locally using
\texttt{model\_training\_script.py}. The application is launched with
\texttt{streamlit run app.py}, which opens the dashboard in the default browser at
\texttt{localhost:8501}.

For cloud-based development, the repository includes a \texttt{.devcontainer}
configuration allowing one-click launch in GitHub Codespaces with all dependencies
pre-installed.

%%─────────────────────────────────────────────────────────────────────────────
\chapter{Future Roadmap}
%%─────────────────────────────────────────────────────────────────────────────

\section{Strategic Vision}

Distillation Column AI is designed as an evolving platform. The current version (v2.2)
delivers a fully functional AI copilot for single-column ethanol purity prediction.
The roadmap extends this foundation towards enterprise-grade process intelligence,
supporting multiple columns, live DCS integration, and advanced time-series analytics.

\begin{table}[H]
\centering
\caption{Development Roadmap}
\label{tab:roadmap}
\renewcommand{\arraystretch}{1.4}
\begin{tabularx}{\textwidth}{|L{3cm}|C{2cm}|L{9cm}|}
\hline
\thead{Version} & \thead{Status} & \thead{Scope and Features} \\
\hline
v1.0 --- ML Core          & Complete  & Physics-informed feature engineering; Random Forest + XGBoost; Streamlit dashboard \\
v2.0 --- AI Intelligence  & Complete  & Gemini diagnostics; SHAP explainability; anomaly detection; smart alerts \\
v2.1 --- AI Chat          & Complete  & Conversational AI copilot; natural-language Q\&A; context-aware multi-turn responses \\
v2.2 --- Optimisation     & Complete  & What-if simulator; sensitivity analysis; target purity solver; energy efficiency advisor \\
v3.0 --- Advanced Analytics & Planned & Time-series trend dashboards; automated shift handover reports; batch comparison \\
v3.1 --- Multi-Column     & Planned   & Support for multiple distillation columns within a single unified dashboard \\
v4.0 --- Edge Deployment  & Future    & Docker container; REST API; OPC-UA connector for live DCS/SCADA data feed \\
\hline
\end{tabularx}
\end{table}

\section{Emerging Technology Integration}

The v4.0 roadmap targets integration with plant-floor automation infrastructure via
OPC-UA and Modbus connectors, transforming the platform from an offline analytical
tool into a real-time process advisor. Potential extensions include fine-tuned Gemini
prompts for domain-specific terminology, LIME as an alternative explainability method,
and neural-network models for capturing non-stationary column dynamics over long
time horizons.

\section{Contribution Areas}

Table~\ref{tab:contrib} documents open areas where community contributions would
be most impactful.

\begin{table}[H]
\centering
\caption{Contribution Areas}
\label{tab:contrib}
\renewcommand{\arraystretch}{1.3}
\begin{tabularx}{\textwidth}{|L{2.5cm}|L{11.5cm}|}
\hline
\thead{Area} & \thead{Suggested Contribution} \\
\hline
AI / LLM         & Fine-tune Gemini prompts for better diagnostic precision; implement GPT-4 as a fallback provider \\
Explainability   & Add LIME as an alternative to SHAP; integrate counterfactual explanation generation \\
Training Data    & Contribute additional real-plant distillation experiments or open-source datasets \\
Models           & Explore neural networks, gradient boosting variants, or ensemble stacking approaches \\
Analytics        & Build time-series dashboards for long-term column performance tracking and drift analysis \\
UI               & Dark mode; mobile-responsive layout; integrated trend charts within the prediction panel \\
Deployment       & Docker image; cloud deployment on AWS, GCP, or Azure; automated CI/CD pipeline \\
Integration      & OPC-UA and Modbus connectors for live DCS/SCADA data ingestion \\
\hline
\end{tabularx}
\end{table}

%%─────────────────────────────────────────────────────────────────────────────
%% BIBLIOGRAPHY
%%─────────────────────────────────────────────────────────────────────────────
\clearpage
\begin{thebibliography}{99}
\addcontentsline{toc}{chapter}{Bibliography}

\bibitem{streamlit}
Streamlit Inc. \textit{Streamlit Documentation}. \url{https://docs.streamlit.io/}

\bibitem{gemini}
Google DeepMind. \textit{Gemini API Documentation}. \url{https://ai.google.dev/}

\bibitem{shap}
S.~M.\ Lundberg and S.-I.\ Lee. A Unified Approach to Interpreting Model Predictions.
\textit{Advances in Neural Information Processing Systems}, 30, 2017.
\url{https://shap.readthedocs.io/}

\bibitem{sklearn}
F.~Pedregosa et~al. Scikit-learn: Machine Learning in Python.
\textit{Journal of Machine Learning Research}, 12:2825--2830, 2011.

\bibitem{xgboost}
T.~Chen and C.~Guestrin. XGBoost: A Scalable Tree Boosting System.
\textit{Proceedings of the 22nd ACM SIGKDD}, 2016.
\url{https://xgboost.readthedocs.io/}

\bibitem{pandas}
The pandas development team. \textit{pandas: Powerful data analysis for Python}.
\url{https://pandas.pydata.org/}

\bibitem{isoforest}
F.~T.\ Liu, K.~M.\ Ting, and Z.-H.\ Zhou. Isolation Forest.
\textit{IEEE International Conference on Data Mining}, 2008.

\bibitem{plotly}
Plotly Technologies Inc. \textit{Collaborative data science}. 2015.
\url{https://plotly.com}

\end{thebibliography}

\end{document}